Bruken av prosedyregenerert innhold har vokst i mange forskjellige bruksområder de siste årene. Dog har ikke prosedyregenert musikk vært like godt dekket som andre domener. Prosedyregenerert musikk har en unik utfordring i den form at det er svært subjektivt hva som kan kalles god musikk.

Denne oppgaven tar for seg velkjente musikalske konsepter og formaliserer dem som \textit{constraints}, altså begrensninger, i et prosedyregenereringssytem. Systemet ble utviklet igjennom tre iterasjoner hvor den siste iterasjonen var assistert med kunstig intelligens for å undersøke hvor nyttig det er som et utviklingsverktøy.

Sluttsystemet kan generere melodier og harmonier som følger teoretiske musikkprinsipper og kan enkelt formes av sluttbruker til å vekte forskjellige musikalske konsepter etter personlig preferanse.
